% ============================================
% 地理参考答案 LaTeX 模板
%  适用于高考地理试卷的答案解析排版
% 使用 XeLaTeX 编译
% ============================================
\documentclass[12pt,a4paper]{ctexart}

% ========== 页面布局 ==========
\usepackage[top=2cm,bottom=2cm,left=2.5cm,right=2.5cm]{geometry}
\usepackage{setspace}
\onehalfspacing

% ========== 列表 ==========
\usepackage{enumitem}

% ========== 数学公式 ==========
\usepackage{amsmath,amssymb}
\usepackage{bm}

% ========== 插图 ==========
\usepackage{graphicx}
\graphicspath{{images/}}

% ========== 颜色 ==========
\usepackage{xcolor}

% ========== 表格 ==========
\usepackage{array,booktabs,caption}

% ========== 页眉页脚 ==========
\usepackage{fancyhdr}
\setlength{\headheight}{13.6pt}
\pagestyle{fancy}
\fancyhf{}
\fancyhead[C]{\small 202X年普通高等学校招生全国统一考试\,$\cdot$\,地理\quad 参考答案}
\fancyfoot[C]{\thepage}
\renewcommand{\headrulewidth}{0.4pt}

% ========== 超链接 ==========
\usepackage[hidelinks]{hyperref}

% ========== 自定义命令 ==========

% 【答案】红色加粗
\newcommand{\daan}[1]{{\color{red}\textbf{【答案】#1}}}

% 【解析】灰色前缀，正文用楷体
\newcommand{\jieti}{\par{\color{gray!60!black}\textbf{【解析】}}}

% 【详解】蓝色前缀
\newcommand{\xijie}{\par{\color{blue!60!black}\textbf{【详解】}}}

% 【小问1详解】、【小问2详解】
\newcommand{\xiaoI}{\par{\color{blue!60!black}\textbf{【小问 1 详解】}}}
\newcommand{\xiaoII}{\par{\color{blue!60!black}\textbf{【小问 2 详解】}}}

% 插入公式图片（WMF 转为 PNG）
\newcommand{\eqimg}[2][0.5]{%
  \includegraphics[width=#1\textwidth,keepaspectratio]{#2}%
}

% 虚数单位和自然底数
\newcommand{\mi}{\mathrm{i}}
\newcommand{\me}{\mathrm{e}}

% ========== 正文 ==========
\begin{document}

% 标题区
\begin{center}
    {\large\textbf{绝密\,$\bigstar$\,启用前\qquad 试卷类型：A}}
    \vspace{1em}
    {\Large\textbf{202X年普通高等学校招生全国统一考试}}
    \vspace{0.5em}
    {\huge\textbf{地\ \ 理}}
    \vspace{1em}
    {\large\textbf{参考答案}}
    \vspace{0.5em}
    \hrule
\end{center}

% ══════════════════════════════════════
%  一、选择题（15 小题，共 45 分）
% ══════════════════════════════════════
\section*{一、选择题}
本题共 15 小题，每小题 3 分，共 45 分。在每小题给出的四个选项中，
只有一项是符合题目要求的。

\begin{enumerate}
    \item 题目文本
    \begin{flushleft}
        A.\ \eqimg[0.15]{imageN.png}\quad
        B.\ \eqimg[0.15]{imageN.png}\quad
        C.\ \eqimg[0.15]{imageN.png}\quad
        D.\ \eqimg[0.15]{imageN.png}
    \end{flushleft}
    \daan{B}
    \jieti 解析内容.
    \xijie 详解内容.
\end{enumerate}

% ══════════════════════════════════════
%  二、非选择题（共 55 分）
% ══════════════════════════════════════
\section*{二、非选择题}
共 55 分。第 16～18 题为必考题，每个试题考生都必须作答。

\begin{enumerate}[resume]
    \item （12分）题目描述．
    \begin{enumerate}
        \item 第1问；
        \item 第2问．
    \end{enumerate}
    \daan{(1) 答案\quad (2) 答案}
    \jieti 解析内容.
    \xiaoI 第1问详解.
    \xiaoII 第2问详解.
\end{enumerate}

\end{document}
